\documentclass[11pt]{article}

\usepackage[spanish,es-tabla]{babel}
\spanishdecimal{.}
\usepackage[T1]{fontenc}
\usepackage[utf8]{inputenc}
\usepackage{newtxtext,newtxmath}
\usepackage{microtype}
\usepackage[margin=2.15cm,headheight=24pt]{geometry}
\usepackage{xcolor}
\usepackage{booktabs}
\usepackage{tabularx}
\usepackage{array}
\usepackage{enumitem}
\usepackage{fancyhdr}
\usepackage{titlesec}
\usepackage{xurl}
\usepackage[hidelinks]{hyperref}

\definecolor{bcp}{HTML}{0B2545}
\definecolor{cobre}{HTML}{C65D12}
\definecolor{tinta}{HTML}{1C2833}

\pagestyle{fancy}
\fancyhf{}
\fancyhead[L]{\textcolor{bcp}{\small MGT209 · E01 · Versión 2}}
\fancyhead[R]{\textcolor{bcp}{\small Banca Personas · BCP}}
\fancyfoot[L]{\textcolor{tinta}{\footnotesize Agente financiero personal}}
\fancyfoot[R]{\textcolor{tinta}{\footnotesize\thepage}}
\renewcommand{\headrulewidth}{0.4pt}
\renewcommand{\headrule}{\hbox to\headwidth{\color{cobre}\leaders\hrule height \headrulewidth\hfill}}
\renewcommand{\footrulewidth}{0pt}

\titleformat{\section}{\large\bfseries\color{bcp}}{\thesection}{0.7em}{}[\vspace{-0.4em}{\color{cobre}\titlerule}]
\titleformat{\subsection}{\normalsize\bfseries\color{bcp}}{\thesubsection}{0.6em}{}
\titlespacing*{\section}{0pt}{1.05em}{0.45em}
\titlespacing*{\subsection}{0pt}{0.75em}{0.25em}

\setlist[itemize]{leftmargin=1.15em,itemsep=0.18em,topsep=0.25em}
\setlist[enumerate]{leftmargin=1.25em,itemsep=0.18em,topsep=0.25em}
\setlength{\parskip}{0.32em}
\setlength{\parindent}{0pt}
\hypersetup{pdftitle={E01 Versión 2. Agente financiero personal},pdfauthor={Franco Stywar Antony Vargas Chino; Albert Christian Hidalgo Sánchez; Sergio Perez Ruiz; Carlos E. Ruiz Perez; Dennis Jack Paitan Cano}}

\newcommand{\usd}[1]{USD~#1}
\newcommand{\pen}[1]{S/~#1}

\begin{document}

\begin{center}
{\color{bcp}\small Universidad Nacional de Ucayali \textperiodcentered\ Maestría en Ingeniería de Sistemas}\\[0.3em]
{\color{cobre}\bfseries\small MGT209 Gestión de Proyectos en TI \textperiodcentered\ Entregable E01 \textperiodcentered\ Versión 2}\\[0.6em]
{\color{bcp}\LARGE\bfseries Caso de negocio y constitución del proyecto}\\[0.3em]
{\large Agente financiero personal de Banca Personas del BCP}\\[0.15em]
{\normalsize Informa y simula. No asesora, no coloca y no mueve dinero.}\\[0.45em]
{\small Franco Stywar Antony Vargas Chino \textperiodcentered\ Albert Christian Hidalgo Sánchez}\\[0.15em]
{\small Sergio Perez Ruiz \textperiodcentered\ Carlos E. Ruiz Perez \textperiodcentered\ Dennis Jack Paitan Cano}\\[0.3em]
{\small Noviembre 2026}
\end{center}

\vspace{0.2em}
{\color{cobre}\rule{\textwidth}{0.6pt}}

El proyecto es un agente financiero personal, disponible en la aplicación y en la web, con capacidades que el cliente concede y puede revocar. Los indicadores salen del Excel \path{E01_Evaluacion_Economica_BCP_Agente.xlsx}. Las celdas amarillas son supuestos del equipo. No hay un VAN publicado de un producto equivalente: se declara así, y la sensibilidad muestra en qué punto el caso deja de cubrir la tasa.

\section{Ficha ejecutiva}

\subsection{Problema}

El cliente de banca personas ya tiene aplicación, banca por internet y Yape. No tiene un agente propio, activo a cualquier hora, que use su posición real y solo las facultades que él le dio. Credicorp informa que BCP Perú cerró 2024 con más de 13,7 millones de clientes y que el 76\,\% opera en digital; Yape tenía 13,7 millones de usuarios activos mensuales, definidos como quienes hicieron al menos una transacción de salida en treinta días (Credicorp, 2025a, 2025b). Esas cifras describen el banco. No son usuarios de este agente.

Lo que ya existe no cubre el hueco. Clara, en WhatsApp, informa saldos, préstamos y fechas de pago (BCP, s.\,f.-a). Salvador orienta a emprendedores y el propio banco dice que no opera, no vende y no aprueba (BCP, s.\,f.-b). En otros mercados el patrón es el mismo con más escala: Erica, de Bank of America, no usa IA generativa ni da consejo de inversión (Bank of America, s.\,f.); AIR, de Revolut, actúa sobre la cuenta del propio cliente y exige biometría para lo sensible (Revolut, 2026); Mago, de Mercado Pago, simula y paga solo con confirmación del titular (Mercado Libre, 2026). Ninguno publica un VAN. En el Perú, Plin por WhatsApp sí ejecuta pagos, con clave y con un tope (Interbank, 2025). Este proyecto no copia esa ejecución.

Los sistemas de los productos no se reemplazan. Cada uno sigue siendo el dueño de sus saldos, obligaciones y movimientos. El agente consulta una vista por cliente, con una marca de si el dato está vigente. Si esa vista no está al día, no contesta con el histórico ni completa el hueco. Esa es la falla de diseño que el proyecto evita.

\subsection{Objetivo}

En diez meses, hasta la salida a producción, ofrecer en Banca Móvil y en la web un agente que informa y simula, con tres perfiles que el cliente concede y revoca. Al cierre del mes 11 el traspaso está hecho. Al tercer año de operación:

\begin{itemize}
  \item 500\,000 clientes mantienen un permiso vigente. Es un supuesto del equipo, cerca del 5\,\% de la base digital si se toma el 76\,\% de 13,7 millones. No es un pronóstico del banco.
  \item El agente no ejecuta ninguna transferencia ni ningún desembolso. Quien confirma es la persona, en el canal; quien ejecuta es el sistema del producto.
  \item Si la vista de posición no está vigente, el agente no afirma un saldo ni lo completa con el histórico.
  \item Yape queda como propuesta al dueño de ese producto. No entra en la meta ni en el beneficio.
\end{itemize}

\subsection{Solución}

Seis habilidades, concedidas una a una: ver la posición, explicar movimientos, simular un crédito del banco, proyectar un plan de gasto contra la deuda, mostrar ofertas que el producto ya tasó, y preparar un encargo para que el canal pida confirmación. Tres perfiles las agrupan: Consulta, Planificación, y Ofertas con confirmación. El detalle está en las secciones 2 y 3.

\subsection{Beneficio que entra al flujo}

Entran tres partidas, todas con supuesto declarado. La primera es el margen del primer año de un crédito que el cliente confirma después de una simulación: en el caso base, el 2\,\% de esas simulaciones. La segunda es la pérdida esperada que se evita cuando la proyección muestra que el plan no cabe y el cliente no toma un crédito que sí habría tomado. La tercera es el margen de una oferta ya tasada, solo si el cliente concedió ese perfil y la acepta en el canal.

No entra el ahorro de Clara. No entra un ingreso por NPS. No entra Yape. No entra asesoría de inversión: recomendar valores corresponde a una sociedad agente autorizada por la SMV, y esa intermediación del grupo no está en esta aplicación (Decreto Supremo N.° 020-2023-EF, 2023).

\section{Cómo se engancha al banco}

El agente consulta lo que el banco ya tiene. No reemplaza los sistemas de cuentas, tarjetas ni créditos.

\begin{itemize}
  \item \textbf{Cada producto sigue siendo el registro.} El saldo y la obligación nacen ahí.
  \item \textbf{Vista vigente.} Por cliente se lee saldo, obligaciones, movimientos recientes, si una oferta ya está calculada y si ese dato está al día. Si no lo está, la herramienta responde que no está confirmado. El texto del modelo no completa el hueco.
  \item \textbf{Histórico.} Sirve para auditar, probar el agente hacia atrás y mejorar. No contesta el saldo de la sesión cuando la vista vigente existe.
  \item \textbf{Herramientas.} Catálogo cerrado. Simular y ejecutar son contratos distintos. Cada llamada revalida el permiso. Sin aprobación no hay efecto, que es la compuerta que FINOS pide para agentes (FINOS, s.\,f.). La Autoridad Monetaria de Singapur coloca una salvaguarda entre el agente y los sistemas: denegar, escalar, ejecutar solo dentro de un umbral, u observar (Monetary Authority of Singapore [MAS], 2026). Aquí el umbral de ejecución es cero.
  \item \textbf{Canales.} La aplicación y la web llaman al agente con la sesión que el cliente ya tiene. No se rediseña la app.
\end{itemize}

Yape no publica un portal de desarrolladores ni una API para que un tercero opere un asistente. Lo público es la aplicación, Yape Empresa y un canal de atención (Yape, s.\,f.). La prensa describe IA en ese canal de consultas; eso no es un punto de enchufe. El proyecto financia un descubrimiento: deep link o una API que el dueño de Yape acepte. No se escribe en sus sistemas. El beneficio del Excel no cuenta un solo usuario de Yape.

Tampoco se hace lo siguiente: inventar un saldo; usar el histórico como si fuera la posición de hoy; desembolsar o mover dinero desde el texto del agente; presentar una simulación como crédito aprobado; ofrecer productos de inversión.

\section{Habilidades y perfiles}

La habilidad es lo que el cliente autoriza. El perfil es el paquete. Revocar corta la herramienta en la llamada siguiente. Ley N.° 29733: consentimiento previo, informado, expreso, para una finalidad determinada, y revocable (Congreso de la República del Perú, 2011).

\begin{table}[ht]
\centering
\small
\begin{tabularx}{\textwidth}{@{}l X X@{}}
\toprule
\textbf{Habilidad} & \textbf{Puede leer o pedir} & \textbf{No puede} \\
\midrule
Ver mi posición & Saldo, obligaciones y frescura & Transferir, ni citar una cifra que la herramienta no devolvió \\
Explicar movimientos & Extracto proyectado & Reclasificar un cargo o abrir una disputa \\
Simular un crédito & Cuota, plazo, TCEA y total, desde el simulador del producto & Desembolsar, reservar línea o llamar a eso una aprobación \\
Proyectar un plan & Obligaciones y una simulación de flujo, por ejemplo un viaje & Comprar, apartar saldo o comprometer un pago \\
Mostrar ofertas & Ofertas ya calculadas por el producto & Aceptarlas o contactar a un tercero \\
Preparar un encargo & Borrador de monto, destino y producto & Enviarlo al producto si la confirmación falta o expiró \\
\bottomrule
\end{tabularx}
\end{table}

\textbf{Consulta} lee posición y movimientos. \textbf{Planificación} suma las dos simulaciones. \textbf{Ofertas con confirmación} suma las ofertas y el borrador. Aceptar, desembolsar y mover dinero los confirma la persona y los ejecuta el producto.

La Resolución SBS N.° 3274-2017, en su texto consolidado al 2 de julio de 2026, exige transparencia de tasas y comisiones, hoja resumen, y un simulador para créditos en cuotas (Superintendencia de Banca, Seguros y AFP [SBS], 2017/2026). La Resolución SBS N.° 01741-2026 exige una vía humana cuando la atención es automatizada, y consentimientos diferenciados, sin casillas premarcadas (SBS, 2026). No se encontró una norma SBS de asesoría automatizada. Por eso el charter no usa la palabra asesor: informa y simula. La cotización es ilustrativa. No es oferta vinculante. Si el cliente contrata, lo hace en el canal, con hoja resumen, TCEA y constancia.

La Resolución SBS N.° 504-2021 pide mínimo privilegio, autenticación y registros de las operaciones que exigen autenticación, y contrato con el proveedor que procese los datos (SBS, 2021). El modelo no se entrena con la conversación del cliente. Si una fase posterior permitiera pagar, entraría la autenticación reforzada. Esa fase no está en este alcance.

\section{Evaluación económica}

El método es el de Sapag Chain (2011): flujo de caja neto, VAN con el año 0 sin descontar, TIR y payback simple. La tasa de 12\,\% es un supuesto del equipo. No es el costo de capital del BCP.

\subsection{Qué está simulado}

No hay un costo por usuario ni un lift auditado para este tipo de agente. Klarna publicó un ahorro estimado de su asistente de servicio, no de un agente de capacidades concedidas (Klarna, 2024). DBS publica el valor de todo su portafolio de modelos, no el VAN de un agente personal (DBS, 2025). Esas cifras no entran al Excel.

El caso base supone, y todo esto está en amarillo: 500\,000 permisos vigentes desde el año 3, con rampas de 16\,\% y 44\,\% en los años 1 y 2; cinco sesiones al año; 30\,\% de ellas son simulación de crédito y otro 30\,\% es proyección de un plan contra la deuda; el 2\,\% de las simulaciones de crédito lo confirma el cliente en el canal; el ticket medio es \pen{6\,000}, unos \usd{1\,600} al tipo de cambio de 3,75, que también es supuesto; el margen del primer año es 6\,\% del ticket y no se repite; en el 35\,\% de las proyecciones el plan no cabe, y el 5\,\% de esos clientes igual habría tomado el crédito, con una pérdida esperada del 3\,\% del ticket; el 30\,\% concede ofertas y el 8\,\% de ellos acepta una al año, con \usd{18} de margen. En el año 3 eso es 15\,000 créditos confirmados y \pen{90} millones de principal. Frente a una cartera de créditos de BCP del orden de \pen{114\,000} millones (Credicorp, 2025a), el principal atribuido es pequeño. Sigue siendo un supuesto, no una meta comercial del banco.

\subsection{Inversión y operación}

El CAPEX base es \usd{1\,420\,000}: vista de posición del cliente, 480\,000; herramientas y consentimiento, 360\,000; adaptadores de la app y la web, 280\,000; descubrimiento con Yape, 80\,000; seguridad, conducta y datos, 140\,000; habilitación, 80\,000. Contingencia de 10\,\% y gestión de 5\,\%, sumadas y no compuestas. La inversión del año 0, y la del charter, es \usd{1\,633\,000}. El caso base supone que la contingencia se gasta.

El OPEX suma un fee de plataforma de \usd{80\,000}, un equipo de \usd{260\,000} y soporte de \usd{40\,000}, estos dos con 3\,\% anual, y \usd{0{,}025} por sesión. El fee no crece con los usuarios.

\begin{table}[ht]
\centering
\small
\begin{tabular}{lrrrrr}
\toprule
 & \textbf{1} & \textbf{2} & \textbf{3} & \textbf{4} & \textbf{5} \\
\midrule
Crédito confirmado & 0{,}23 & 0{,}63 & 1{,}44 & 1{,}44 & 1{,}44 \\
Pérdida evitada & 0{,}10 & 0{,}28 & 0{,}63 & 0{,}63 & 0{,}63 \\
Ofertas & 0{,}03 & 0{,}10 & 0{,}22 & 0{,}22 & 0{,}22 \\
OPEX & 0{,}39 & 0{,}42 & 0{,}46 & 0{,}47 & 0{,}48 \\
\bottomrule
\end{tabular}
\end{table}

\noindent Millones de USD. El año 0 es solo la inversión, \usd{1{,}63} millones, sin beneficio ni OPEX.

\begin{table}[ht]
\centering
\small
\begin{tabular}{lrrrrrr}
\toprule
 & \textbf{0} & \textbf{1} & \textbf{2} & \textbf{3} & \textbf{4} & \textbf{5} \\
\midrule
Flujo de caja neto & $-1{,}63$ & $-0{,}02$ & 0{,}59 & 1{,}83 & 1{,}82 & 1{,}81 \\
\bottomrule
\end{tabular}
\end{table}

\begin{table}[ht]
\centering
\small
\begin{tabularx}{\textwidth}{@{}l l X@{}}
\toprule
\textbf{Indicador} & \textbf{Caso base} & \textbf{Lectura} \\
\midrule
VAN al 12\,\% & \usd{2\,292\,923} & Cubre la tasa del equipo. \\
TIR & 43{,}0\,\% & La tasa de descuento no está dentro de la TIR. \\
Payback simple & 2{,}59 años & El acumulado cruza cero durante el año 3 de operación. \\
Beneficios & \usd{8\,229\,600} & Cinco años, sin descontar. \\
OPEX & \usd{2\,217\,741} & Fee, equipo, soporte e inferencia. \\
TCO & \usd{3\,850\,741} & Inversión más OPEX. \\
ROI sobre la inversión & 268{,}1\,\% & (Beneficios $-$ TCO) / inversión. El CAPEX no se resta dos veces. \\
\bottomrule
\end{tabularx}
\end{table}

El primer año de operación todavía no paga su OPEX. El resultado se arma desde el segundo año, cuando los permisos vigentes pasan de 80\,000 a 220\,000 y después a 500\,000.

\subsection{Dónde se cae}

\begin{table}[ht]
\centering
\small
\begin{tabular}{lrrr}
\toprule
\textbf{Escenario} & \textbf{VAN} & \textbf{TIR} & \textbf{ROI} \\
\midrule
Base & \usd{2\,292\,923} & 43{,}0\,\% & 268{,}1\,\% \\
Conversión 1{,}5\,\% & \usd{1\,425\,918} & 32{,}7\,\% & 188{,}8\,\% \\
Conversión 1{,}0\,\% & \usd{558\,913} & 20{,}8\,\% & 109{,}4\,\% \\
Conversión 0{,}8\,\% & \usd{212\,111} & 15{,}5\,\% & 77{,}7\,\% \\
Sin pérdida evitada & \usd{775\,664} & 24{,}0\,\% & 129{,}3\,\% \\
Meta de 250\,000 clientes & \usd{$-384\,557$} & 5{,}1\,\% & 23{,}1\,\% \\
\bottomrule
\end{tabular}
\end{table}

Con la pérdida evitada, la conversión puede bajar hasta cerca del 0{,}7\,\% antes de igualar el 12\,\%. Si esa partida se apaga, porque alguien no acepta monetizar un incumplimiento que no ocurrió, el caso todavía pasa: la TIR queda en 24{,}0\,\%, y el umbral de conversión sube a cerca del 1{,}6\,\%. Lo que sí tumba el caso es la adopción. Con 250\,000 permisos vigentes la TIR cae a 5{,}1\,\%. La defensa no es «el modelo rinde solo». Es «rinde si el cliente concede el permiso y confirma en el canal».

\section{Project charter}

\textbf{Nombre.} Agente financiero personal de Banca Personas.

\textbf{Propósito.} Que el cliente tenga un agente 24/7, en la aplicación y en la web, limitado a lo que él autorizó y apoyado en la posición vigente que ya guardan los productos, no en un texto libre ni en el histórico.

\textbf{Patrocinio.} Un solo dueño del beneficio: la gerencia de banca personas digital, responsable de los permisos vigentes y de que el cliente confirme en el canal. El CIO patrocina el enganche a los sistemas que ya existen; no comparte la meta de 500\,000 permisos. Finanzas aprueba el desembolso de \usd{1\,633\,000} y la contingencia. El director del proyecto puede frenar una salida que afirme un saldo sin vista vigente o que presente una simulación como crédito aprobado. No puede prometer la conversión del 2\,\%.

\textbf{Alcance.} Vista de posición para cuentas, tarjetas y créditos de personas; herramientas y consentimiento; tres perfiles; adaptadores de la app y la web; vía humana; descubrimiento con Yape, sin construirlo.

\textbf{Fuera de alcance.} Desembolsos, transferencias y pagos ejecutados por el agente. Asesoría de inversión o de crédito. Reemplazo de Clara o de Salvador. Rediseño de la aplicación, de la web o de Yape. Que el agente escriba por su cuenta en los sistemas de los productos. Beneficio económico atribuido a usuarios de Yape. Entrenar el modelo con conversaciones de clientes.

\textbf{Hitos.} Mes 2: arquitectura, amenaza y diseño del consentimiento. Mes 4: vista de cuentas, tarjetas y créditos en ambiente de prueba, sin modificar esos sistemas. Mes 6: piloto del perfil Consulta. Mes 8: Planificación en la app y en la web. Mes 10: Ofertas con confirmación en producción, con escape a una persona. Mes 11: cierre, traspaso y decisión registrada sobre Yape. El mes 11 no es ejecución. Los beneficios se miden en años de operación.

\textbf{Presupuesto.} \usd{1\,633\,000} en el año 0. Es la misma cifra del flujo.

\textbf{Riesgos que el charter asume.} Vista vencida contestada con una cifra. Simulación leída como aprobación. Alcance que se corre hacia Yape para hacer cuadrar el VAN. Permisos vigentes por debajo de 250\,000. Conversión por debajo del umbral de la sección 4.

\section{Matriz de interesados}

El criterio es el de poder, interés e impacto de la \textit{Guía de los fundamentos para la dirección de proyectos} (Project Management Institute, 2021).

\begin{table}[ht]
\centering
\small
\begin{tabularx}{\textwidth}{@{}l l X@{}}
\toprule
\textbf{Interesado} & \textbf{Estrategia} & \textbf{Por qué} \\
\midrule
Banca personas digital & Gestionar de cerca & Único dueño del beneficio: permisos y confirmación del cliente. \\
CIO y plataforma de datos & Gestionar de cerca & Dueños de los sistemas de los productos. El proyecto solo se engancha. \\
Riesgo de crédito & Gestionar de cerca & La pérdida evitada y el simulador son de su dominio. \\
Oficial de conducta de mercado & Gestionar de cerca & Frontera entre informar y asesorar. Valida el diseño (SBS, 2017/2026). \\
CISO & Mantener satisfecho & Mínimo privilegio, trazas y contrato con el proveedor del modelo. \\
Finanzas & Mantener satisfecho & Desembolso y contingencia. No define las habilidades. \\
Dueño de Yape & Mantener satisfecho & Poder alto e interés bajo en este alcance. No es sponsor. \\
Dueño de Clara & Involucrar & Para no solapar el canal de WhatsApp ni contar dos veces el mismo ahorro. \\
Clientes & Involucrar & Poder formal bajo. Sin su permiso no hay sesión ni beneficio. \\
SBS & Cumplir & Conducta, simulador y vía humana. No se negocia el VAN. \\
SMV y la sociedad agente del grupo & Quedarse fuera & Valores y asesoría de inversión no entran. \\
Proveedor del modelo & Gestionar de cerca & Disponibilidad, y la prohibición de entrenar con datos del cliente. \\
\bottomrule
\end{tabularx}
\end{table}

\section{Cierre}

La inversión de \usd{1{,}63} millones se justifica en el caso base: VAN de \usd{2{,}29} millones, TIR de 43{,}0\,\% y recuperación en 2{,}59 años. También se justifica si no se cuenta la pérdida evitada. No se justifica si los permisos vigentes se quedan en la mitad de la meta. El agente informa y simula sobre la posición vigente del cliente. El cliente concede cada habilidad. El producto, y no el modelo, ejecuta. Yape se propone; no se presupone.

\begin{thebibliography}{99}
\bibitem[BCP(s.\,f.-a)]{bcp-clara}
Banco de Crédito del Perú. (s.\,f.-a). \textit{Consúltale a Clara}. \url{https://www.viabcp.com/canales/preguntale-a-clara}

\bibitem[BCP(s.\,f.-b)]{bcp-salvador}
Banco de Crédito del Perú. (s.\,f.-b). \textit{Asesor IA de negocio}. \url{https://www.viabcp.com/contigo-emprendedor-bcp/asesor-ia-ce-whatsapp}

\bibitem[BofA(s.\,f.)]{erica}
Bank of America. (s.\,f.). \textit{Erica, your virtual financial assistant}. \url{https://info.bankofamerica.com/en/digital-banking/erica}

\bibitem[Congreso(2011)]{ley29733}
Congreso de la República del Perú. (2011). \textit{Ley N.° 29733, Ley de Protección de Datos Personales}. \url{https://www.leyes.congreso.gob.pe/Documentos/Leyes/29733.pdf}

\bibitem[Credicorp(2025a)]{credicorp-ar}
Credicorp Ltd. (2025a). \textit{Annual and sustainability report 2024}. \url{https://credicorp.gcs-web.com/system/files-encrypted/nasdaq_kms/assets/2025/02/28/17-34-02/Annual%20and%20Sustainability%20Report%202024%20-%20ENG_.pdf}

\bibitem[Credicorp(2025b)]{credicorp-6k}
Credicorp Ltd. (2025b). \textit{Results for the quarter ended December 31, 2024} (Form 6-K, Exhibit 99.1). \url{https://www.sec.gov/Archives/edgar/data/1001290/000114036125004694/ef20043470_ex99-1.htm}

\bibitem[DBS(2025)]{dbs}
DBS. (2025). \textit{Annual report 2024: Innovating impactful solutions}. \url{https://www.dbs.com/annualreports/2024/innovating-impactful-solutions-for-customers.html}

\bibitem[DS(2023)]{smv}
Ministerio de Economía y Finanzas. (2023). \textit{Decreto Supremo N.° 020-2023-EF, Texto Único Ordenado de la Ley del Mercado de Valores}. \url{https://www.smv.gob.pe/ServicioConsultaNormas/uploads/TUO_LMV_DECRETO_SUPREMO_N_020_2023_EF.pdf}

\bibitem[FINOS(s.\,f.)]{finos}
FINOS. (s.\,f.). \textit{MI-24: Human-in-the-loop action approval gate}. \url{https://air-governance-framework.finos.org/mitigations/mi-24_human-in-the-loop-action-approval-gate.html}

\bibitem[Interbank(2025)]{plin}
Interbank. (2025, 31 de octubre). \textit{Plin WhatsApp}. \url{https://interbank.pe/blog/sala-de-prensa/plin-whatsapp-interbank-revolucion-pagos-peru}

\bibitem[Klarna(2024)]{klarna}
Klarna. (2024, 27 de febrero). \textit{AI assistant handles two-thirds of customer service chats in its first month}. \url{https://www.klarna.com/international/press/klarna-ai-assistant-handles-two-thirds-of-customer-service-chats-in-its-first-month/}

\bibitem[MAS(2026)]{safr}
Monetary Authority of Singapore. (2026). \textit{Safeguards for agentic finance at runtime} (SAFR v1.1). \url{https://www.mas.gov.sg/-/media/mas-media-library/development/fintech/ai-safr/safr-v1-1.pdf}

\bibitem[Mercado Libre(2026)]{mago}
Mercado Libre. (2026, 13 de agosto). \textit{Mercado Pago lanza Mago}. \url{https://news.mercadolivre.com/pt/mercado-pago-lanca-mago-assistente-financeiro}

\bibitem[PMI(2021)]{pmbok}
Project Management Institute. (2021). \textit{Guía de los fundamentos para la dirección de proyectos} (Guía del PMBOK) (7.ª ed.).

\bibitem[Revolut(2026)]{revolut}
Revolut. (2026, 9 de abril). \textit{AIR: AI by Revolut}. \url{https://www.revolut.com/blog/post/air-ai-by-revolut/}

\bibitem[Sapag(2011)]{sapag}
Sapag Chain, N. (2011). \textit{Proyectos de inversión: Formulación y evaluación} (2.ª ed.). Pearson Educación.

\bibitem[SBS(2021)]{sbs504}
Superintendencia de Banca, Seguros y AFP. (2021). \textit{Resolución SBS N.° 504-2021, Reglamento para la gestión de la seguridad de la información y la ciberseguridad}. \url{https://busquedas.elperuano.pe/dispositivo/NL/1929393-1}

\bibitem[SBS(2017/2026)]{sbs3274}
Superintendencia de Banca, Seguros y AFP. (2017/2026). \textit{Resolución SBS N.° 3274-2017, Reglamento de gestión de conducta de mercado} (texto consolidado al 2 de julio de 2026). \url{https://intranet2.sbs.gob.pe/dv_int_cn/1731/v12.0/Adjuntos/3274-2017%20R%2002.07.26.pdf}

\bibitem[SBS(2026)]{sbs1741}
Superintendencia de Banca, Seguros y AFP. (2026). \textit{Resolución SBS N.° 01741-2026}. \url{https://intranet2.sbs.gob.pe/dv_int_cn/2589/v1.0/Adjuntos/R%2001741-2026.pdf}

\bibitem[Yape(s.\,f.)]{yape}
Yape. (s.\,f.). \textit{Yape Empresa}. \url{https://www.yape.com.pe/productos/yape-empresa}
\end{thebibliography}

\end{document}
